\section{Zadanie 5}
\autor{Kacper Orszulak}
Niech $T_n$ będzie pełnym drzewem binarnym o głębokości $n$ (czyli $T_n$ ma $2^n - 1$ wierzchołków). Każda krawędź $T_n$ zostaje usunięta z prawdopodobieństwem $1/2$ i niezależnie od akcji na innych krawędziach. Jaka jest oczekiwana liczba wierzchołków w spójnej składowej z korzeniem w wynikowym lesie?

\subsection*{Rozwiązanie}
Jakie jest prawdopodobieństwo, że wierzchołek $v$ na głębokości $k$ będzie w spójnej składowej korzenia? $\mathcal{P}(I_v = 1) = \left(\frac{1}{2}\right)^k$, ponieważ jednocześnie wszystkie $k$ krawędzi od korzenia do $v$ muszą być nieusunięte. Zatem $\mathbb{E}(I_k) = 0 \cdot \mathcal{P}(I_v = 0) + 1 \cdot \mathcal{P}(I_v = 1) = \mathcal{P}(I_v = 1)$. Niech $X = \sum_{v \in V} I_v$ to zmienna losowa liczby wierzchołków w spójnej składowej korzenia. Dostajemy
\[
    \mathbb{E}(X) = \mathbb{E}\left(\sum_{v \in V} I_v\right) = \sum_{v \in V} \mathbb{E}(I_v) = \sum_{v \in V} \mathcal{P}(I_v = 1) = \sum_{v \in V} \left(\frac{1}{2}\right)^k = \sum_{k=0}^{n-1} 2^k \cdot \left(\frac{1}{2}\right)^k = n.
\]

\subsection*{Odpowiedź}
\[
\boxed{n}
\]
